\begin{ledgroupsized}[r]{120mm}
\footnotesize
\pstart
\noindent\textbf{Überlieferung:}
\pend
\end{ledgroupsized}
\begin{ledgroupsized}[r]{114mm}
\footnotesize
\pstart
\parindent -6mm
\makebox[6mm][l]{\textit{L}}%
Aufzeichnung:
LH~XXXV~9, 10~Bl.~7. 
Ein Blatt (ca.~16,5~x~14~cm.);		%
\href{\Scan{LH_35_09_10_007wz}}%
{Wasserzeichenfragment} am Rand (Papier aus Relliehausen bei Hannover);
alle Ränder beschnitten.
Eine Seite auf Bl.~7~r\textsuperscript{o}; Bl.~7~v\textsuperscript{o} leer.
\pend
\end{ledgroupsized}
\vspace{5mm}
\begin{ledgroup}
\footnotesize
\pstart
\noindent
\textbf{Datierungsgründe:}
Die Datierung der Aufzeichnung N.\;\ref{41137} beruht auf dem Wasserzeichen, das im Leibniz-Korpus für den Zeitraum von April 1677 bis Oktober 1684 häufig belegt ist.
\pend
\pstart
Leibniz ist hier um einen Beweis des Satzes bemüht, dass jede flache Figur einen Schwerpunkt besitzt. Damit weist die Aufzeichnung mit dem eigenhändig auf März 1678 datierten Konzept "De centro gravitatis" (N.\;\ref{60279}) nicht nur eine thematische Ähnlichkeit, sondern auch viele inhaltliche Übereinstimmungen auf.
Allerdings sind die festgestellten Anklänge nicht derart eng (etwa wörtliche Übereinstimmung), um eindeutig auf ein Abhängigkeitsverhältnis zwischen N.\;\ref{41137} und dem Konzept N.\;\ref{60279} schließen zu lassen und eine Präzisierung der Datierung auf Grundlage des Wasserzeichens zu ermöglichen.
\pend
\end{ledgroup}
\selectlanguage{latin}
\vspace{8mm}
\pstart
\centering
\href{\Scan{LH_35_09_10_007r}}%
{\lbrack7~r\textsuperscript{o}\rbrack}
\hfill
Demonstratio\indexS{demonstratio centri gravitatis} quod in quolibet corpore sit Centrum Gravitatis\indexS{centrum gravitatis}
\hfill
\phantom{\lbrack7~r\textsuperscript{o}\rbrack}
\pend
\nobreak
\vspace{0.5\baselineskip}
\pstart\noindent 
Centrum gravitatis\indexS{centrum gravitatis} voco, cui si insistat corpus\indexS{corpus} in neutram partem inclinet. 
Ostendam primum id
\edtext{esse in qualibet superficie.}{\lemma{esse in}\Bfootnote{%
\textit{(1)}~quolibet plano %
\textit{(2)}~qualibet superficie.~\textit{L}}}
Nam in qualibet linea esse manifestum est. 
\edtext{In superficie}{\lemma{In}\Bfootnote{%
\textit{(1)}~plano %
\textit{(2)}~superficie~\textit{L}}}
autem ostendendum omnes axes aequilibrii\indexS{axis aequilibrii} se in 
\edtext{eodem puncto}{%
\lemma{eodem}%
\Bfootnote{%
\textit{(1)}~pl
\textit{(2)}~\bars{se \textit{streicht Hrsg.}}
\textit{(3)}~puncto~\textit{L}}}
secare.
\pend
\begin{edfigure}%
\includegraphics[scale=1]{%
gesamttex/edit_VIII-4/images/LH_35_09_10_007_d.pdf}
\edcaption{\lbrack\textit{Fig.~1}\rbrack}
\end{edfigure}
\pstart
Sit superficies \textit{AB}. 
Sit aliquis axis aequilibrii\indexS{axis aequilibrii} qui utique necessario est 
(eodem modo id probatur, quo in puncto probatur esse centrum\indexS{centrum gravitatis} 
\edtext{gravitatis) nam}{\lemma{gravitatis)}\Bfootnote{%
\textit{(1)}~cum enim 
\textit{(2)}~nunc %
\textit{(3)}~nam~\textit{L}}}
corpus semper parallelum sibi huc illuc movendo super recta \textit{DE},
ita scilicet, ut eaedem semper maneant rectae in corpore eidem axi\indexS{axis aequilibrii} vel jugo \textit{DE} 
\edtext{parallelae, patet}{\lemma{parallelae,}\Bfootnote{%
\textit{(1)}~tandem %
\textit{(2)}~patet~\textit{L}}}
corpus nunc in hoc nunc in illud latus praeponderaturum et ideo necessario 
in medio alicubi quieturum, quod ostendi posset operosius. 
Sit ille axis\indexS{axis aequilibrii} \textit{FG}. Diminuatur ille utrobique demtis
\edtext{\lbrack\textit{FH}, \textit{LG}\rbrack}{\lemma{}\Bfootnote{\textit{HL} \textit{L ändert Hrsg.}}}
tali modo ut incumbat corpus \textit{AB} soli axi \textit{HL}.
Nihilominus patet ab una parte adhuc tantundem ponderare quantum ab alia, nec magis in \textit{B} 
quam 
\edtext{in \textit{A} inclinaturum}{\lemma{}%
\Bfootnote{in \textit{A} \bars{in \textit{streicht Hrsg.}} inclinaturum~\textit{L}}}
(nec enim nunc adhuc disputo an ad \textit{F} vel \textit{G} inclinaturum sit).
Cumque \textit{HL} possit esse quantulacunque, poterit esse infinite parva seu punctum ut \textit{M}. 
Jam punctum hoc non magis respectum habet ad hanc rectam
\edtext{\textit{FG}}{\lemma{}\Bfootnote{\textit{FG} \textit{erg.~L}}}
quam ad aliam quamcunque. %
\pend
\pstart
Sed ut inde distinctius consequentiam ducam\lbrack,\rbrack\ invenimus punctum \textit{M}, 
in quo positum corpus non magis inclinat versus \textit{B}, quam versus \textit{A}. 
Ergo etiam non magis versus \textit{N}, quam versus \textit{P}. Nam si inclinat versus \textit{N}, 
magis inclinabit versus \textit{B} quam versus \textit{A}.
Eodem modo si premat in \textit{F} potius quam in \textit{G} magis premet versus \textit{N} quam versus \textit{P}. 
Et ita de qualibet recta, ergo in nullam magis partem premet, 
adeoque unum erit punctum commune aequilibrii.\indexS{punctum commune aequilibrii}
Sed vellem ex hoc solo quod in punctum definitur rem posse duci,
quod punctum illud omnibus axibus\indexS{axis aequilibrii} commune\indexS{punctum omnibus axibus aequilibrii commune} est. 
Nimirum, si punctum hoc \textit{M} datum sit non vero axis,\indexS{axis aequilibrii}
non poterit inde per regressum inveniri ex quo axe sit ortum. 
Verum hoc clariori explicatione indiget.
\pend
